%% kapa: exam=1 round=25 year=2014 period=2 form=A subject=appraisal_law subject_ko=부동산관계법규 seq=16 number=16
%% source: Q-Net 감정평가사 기출문제 (cst00309) artlSeq=5101207 file=extracted/1008/5101207/2014년 제25회 감정평가사 1차시험 기출문제/2교시(A형).hwp
%% conversion: kordoc 4.24.51 -> tools/gen1.py (auto); review=proofread-sonnet; flags=none
\begin{question}{16}
국토의 계획 및 이용에 관한 법령상 용도구역에 관한 설명으로 옳은 것은?
\begin{choices}
  \item 도시의 자연환경 및 경관을 보호하고 도시민에게 건전한 여가ㆍ휴식공간을 제공하기 위하여 도시지역 안에서 식생이 양호한 산지의 개발을 제한할 필요가 있다고 인정되는 지역을 개발제한구역으로 지정할 수 있다.
  \item 국방부장관의 요청에 따라 국토교통부장관이 개발제한구역을 지정하는 경우에는 이를 광역도시계획으로 결정한다.
  \item 국토교통부장관은 도시자연공원구역의 지정을 도시ㆍ군관리계획으로 결정할 수 있다.
  \item 시가화조정구역의 시가화 유보기간은 10년 이상 20년 이내이다.
  \item 시가화조정구역의 지정에 관한 도시ㆍ군관리계획의 결정은 시가화 유보기간이 끝난 날의 다음날부터 그 효력을 잃는다.
\end{choices}
\end{question}
\answer{5}
